%! The Inverse and Determinant of a Matrix — Unit 4, Lesson 4.11 — Homework
\documentclass[10pt]{article}
\usepackage{precalculus-article}
\usepackage{precalculus-boxes}
\newcommand{\TallMath}[1]{$\displaystyle #1\rule[-1.4em]{0pt}{3.2em}$}

\begin{document}

\pageheader{Unit 4, Lesson 4.11}{Homework}
\namedateperiod

\vspace{0.1in}

\begin{notesbox}{Practice}
\begin{enumerate}[leftmargin=1.5em, itemsep=6pt]

    \item Compute the determinant of each matrix.
    \begin{enumerate}[label=(\alph*), leftmargin=1.5em, itemsep=10pt]
        \item $\begin{bmatrix}4&1\\2&3\end{bmatrix}$ \hfill $\det = $ \blank{3cm}

        \item $\begin{bmatrix}-2&5\\1&-3\end{bmatrix}$ \hfill $\det = $ \blank{3cm}

        \item $\begin{bmatrix}6&2\\3&1\end{bmatrix}$ \hfill $\det = $ \blank{3cm}
    \end{enumerate}
    \vspace{0.1in}

    \item For each matrix in problem 1, determine whether it is invertible or singular.
          Show your reasoning.
          \vspace{0.5in}

    \item Find the inverse of each invertible matrix from problem 1.
          Show all work using the formula $A^{-1} = \dfrac{1}{\det(A)}\begin{bmatrix}d&-b\\-c&a\end{bmatrix}$.
          \vspace{1.4in}

    \item Verify your answer to 3(a) by computing $A \cdot A^{-1}$ and confirming the result
          equals $I_2$.
          \vspace{1.0in}

    \item The columns of $M = \begin{bmatrix}4&2\\6&3\end{bmatrix}$ form two vectors.
    \begin{enumerate}[label=(\alph*), leftmargin=1.5em, itemsep=10pt]
        \item Compute $\det(M)$.
              \vspace{0.4in}
        \item Interpret $|\det(M)|$ geometrically: what does it tell you about the column vectors?
              \writeline
              \writeline
    \end{enumerate}

    \item \textbf{AP-Style Multiple Choice.} Which of the following matrices does NOT have an
          inverse?
    \begin{enumerate}[label=(\Alph*), leftmargin=2em, itemsep=4pt]
        \item $\begin{bmatrix}1&0\\0&1\end{bmatrix}$
        \item $\begin{bmatrix}3&1\\6&2\end{bmatrix}$
        \item $\begin{bmatrix}2&-1\\3&4\end{bmatrix}$
        \item $\begin{bmatrix}5&2\\1&1\end{bmatrix}$
    \end{enumerate}

\end{enumerate}
\end{notesbox}

\vspace{0.1in}

\begin{extensionbox}
    Let $A = \begin{bmatrix}a&0\\0&d\end{bmatrix}$ be a diagonal matrix. Derive a formula
    for $A^{-1}$ (assuming $a \ne 0$ and $d \ne 0$) and explain what it means geometrically.
\end{extensionbox}

\end{document}
